\clearpage

\begin{center}
    {\Large \textbf{Anexo: Tabela de Referência — Funções e Métodos Python}}
\end{center}

\vspace{0.15cm}

%\noindent A tabela a seguir sintetiza a sintaxe e a funcionalidade dos métodos e funções nativas (\emph{built-ins}) da linguagem Python. Os exemplos apresentados são intencionalmente genéricos e descontextualizados, servindo como guia rápido de consulta de sintaxe.

\vspace{0.3cm}

\begin{center}
\renewcommand{\arraystretch}{1.22}
\small
\begin{tabularx}{\textwidth}{@{} >{\ttfamily}l l >{\RaggedRight}X >{\ttfamily}l @{}}
\toprule
\textbf{\textsf{Função / Método}} & \textbf{\textsf{Tipo}} & \textbf{\textsf{Descrição / Funcionalidade}} & \textbf{\textsf{Exemplo de Uso}} \\
\midrule
.append(item) & Método (\texttt{list}) & Insere \texttt{item} ao final da lista original (\emph{in-place}). & itens.append("novo") \\
.copy() & Método (\texttt{list}) & Cria cópia rasa (\emph{shallow copy}) independente da lista. & nova = dados.copy() \\
.split([sep]) & Método (\texttt{str}) & Divide a string por \texttt{sep} e retorna lista de substrings. & partes = s.split(",") \\
.strip() & Método (\texttt{str}) & Remove espaços em branco do início e final da string. & limpo = s.strip() \\
.upper() & Método (\texttt{str}) & Converte caracteres alfabéticos da string para maiúsculas. & maiusc = s.upper() \\
.join(seq) & Método (\texttt{str}) & Concatena elementos de \texttt{seq} usando o delimitador base. & " - ".join(palavras) \\
.get(k[, d]) & Método (\texttt{dict}) & Retorna o valor da chave \texttt{k}, ou padrão \texttt{d} se inexistente. & val = mapa.get("k", 0) \\
\midrule
len(seq) & \emph{Built-in} & Retorna a quantidade de elementos de uma coleção. & tam = len(dados) \\
range(i, f[, p]) & \emph{Built-in} & Gera inteiros de \texttt{i} até antes de \texttt{f} com passo \texttt{p}. & range(1, 10, 2) \\
round(num[, n]) & \emph{Built-in} & Arredonda número decimal para $n$ casas decimais (padrão $0$). & res = round(total, 2) \\
sum(iteravel) & \emph{Built-in} & Retorna o somatório dos elementos numéricos do iterável. & soma = sum(valores) \\
int(x) & \emph{Built-in} & Converte número ou string para inteiro truncando decimais. & num = int(4.9) \\
min(a, b, ...) & \emph{Built-in} & Retorna o menor valor entre os argumentos fornecidos. & menor = min(limite, x) \\
max(a, b, ...) & \emph{Built-in} & Retorna o maior valor entre os argumentos fornecidos. & maior = max(base, y) \\
super() & \emph{Built-in} & Acessa métodos e construtores da superclasse na herança. & super().\_\_init\_\_(arg) \\
print(*args) & \emph{Built-in} & Imprime argumentos formatados no console padrão de saída. & print("Total:", val) \\
\bottomrule
\end{tabularx}
\end{center}
